%=======================================================================
% main.tex  --  IEEEtran (conference)
%
% Artigo: degradacao de classificadores CNN por pre-processamento
%         perceptualmente benefico.
%
% Estrutura da Secao II alinhada a rubrica de avaliacao:
%   QP (1,5) II.A | SB (1,0) II.B | CI (1,0) II.C | CE (1,0) II.D
%   FB (1,0) II.E | RR (2,5) II.F | RQ (2,0) II.G
%=======================================================================

\documentclass[conference]{IEEEtran}
\IEEEoverridecommandlockouts

\usepackage[T1]{fontenc}
\usepackage[utf8]{inputenc}
% Se o Overleaf reclamar "Unknown option 'brazil'", basta comentar a linha
% abaixo: o documento compila normalmente sem ela (perde-se apenas a
% hifenizacao em portugues e os nomes traduzidos de "Tabela"/"Figura").
\usepackage[brazil]{babel}
\usepackage{cite}
\usepackage{amsmath,amssymb,amsfonts}
\usepackage{algorithmic}
\usepackage{graphicx}
\usepackage{textcomp}
\usepackage{xcolor}
\usepackage{tikz}
\usetikzlibrary{shapes.geometric, arrows.meta, positioning}

\def\BibTeX{{\rm B\kern-.05em{\sc i\kern-.025em b}\kern-.08em
    T\kern-.1667em\lower.7ex\hbox{E}\kern-.125emX}}

\begin{document}

\title{Quando Melhorar Piora: Uma Revisão Sistemática sobre a Degradação de CNNs por Pré-Processamento de Imagens}

\author{\IEEEauthorblockN{1\textsuperscript{st} Beatriz Almeida Felicio}
\IEEEauthorblockA{\textit{Instituto de Informática} \\
\textit{Universidade Federal de Goiás}\\
Goiânia, Brasil \\
beatrizfelicio@discente.ufg.br}
\and
\IEEEauthorblockN{2\textsuperscript{nd} Júlio Cesar Alves Pereira}
\IEEEauthorblockA{\textit{Instituto de Informática} \\
\textit{Universidade Federal de Goiás}\\
Goiânia, Brasil \\
julio\_alves@discente.ufg.br}
\and
\IEEEauthorblockN{3\textsuperscript{rd} Kamilla Jordana R. de Moura}
\IEEEauthorblockA{\textit{Instituto de Informática} \\
\textit{Universidade Federal de Goiás}\\
Goiânia, Brasil \\
kamillaregosino@discente.ufg.br}
}

\maketitle

\begin{abstract}
Técnicas clássicas de pré-processamento de imagem, como redução de ruído,
suavização, equalização de histograma, aguçamento e super-resolução, são
rotineiramente aplicadas a montante de classificadores baseados em redes
neurais convolucionais sob a suposição de que uma imagem visualmente melhor
é também uma entrada melhor para o modelo. Este trabalho investiga se essa
suposição se sustenta, partindo da hipótese de que tais transformações alteram
ou removem componentes espectrais, particularmente em altas frequências, que
as redes exploram como sinais discriminativos. Apresenta-se uma revisão
sistemática da literatura conduzida para fundamentar a investigação, cobrindo o
protocolo de busca, a seleção dos estudos, a síntese crítica das evidências e
as respostas consolidadas às questões de pesquisa formuladas. A síntese de 29
estudos indica que a degradação é consistente e monotônica para suavização e
substancial para redução de ruído aplicada a imagens já limpas. Além disso, a degradação 
não é uniforme entre técnicas: a equalização adaptativa de histograma melhora o
desempenho sob o mesmo protocolo em que a redução de ruído o piora. Quatro
mecanismos complementares são identificados na literatura. Constatou-se também que a relação entre a
energia espectral removida e a queda de acurácia revela-se apenas parcialmente
preditiva, sendo fortemente mediada pela distribuição espectral do conjunto de treino e pela
simetria do protocolo entre treino e teste.
\end{abstract}

\begin{IEEEkeywords}
pré-processamento de imagem, redes neurais convolucionais, análise espectral,
robustez, classificação de imagens, revisão sistemática
\end{IEEEkeywords}

%=======================================================================
\section{Introdução}
%=======================================================================

\textit{Pipelines} de visão computacional frequentemente inserem uma etapa de
pré-processamento entre a aquisição da imagem e o classificador. A motivação
declarada costuma ser a melhoria da qualidade da entrada: remover ruído,
atenuar artefatos de compressão, realçar bordas, equalizar contraste ou
aumentar a resolução. A suposição implícita é que o classificador se beneficia
das mesmas transformações que tornam a imagem mais agradável ao observador
humano.

Essa suposição raramente é testada. Quando o é, o resultado tende a
contrariá-la. Há evidência acumulada de que classificadores convolucionais
exploram, como sinal de decisão, componentes da imagem que o sistema visual
humano descarta --- textura local, detalhe fino e energia espectral em bandas
que o olho não resolve. Se isso é verdade, então as técnicas de aprimoramento
perceptual atuam precisamente sobre os elementos que o modelo usa para decidir.
Consequentemente, o ganho de aparência é pago com a perda de acurácia. A Fig.~\ref{fig:hipotese_conceitual}
sintetiza esse compromisso.

A questão central deste trabalho é investigar se técnicas clássicas de
pré-processamento e aprimoramento perceptual de imagens --- incluindo redução
de ruído, suavização, equalização de histograma, aguçamento e super-resolução
--- podem degradar sistematicamente o desempenho de classificadores CNN
pré-treinados, mesmo produzindo imagens visualmente mais agradáveis para
observadores humanos. A hipótese subjacente é que essas transformações
modificam ou removem componentes espectrais, particularmente em altas
frequências, que são explorados pelas CNNs como sinais discriminativos
relevantes para a classificação. Além disso, busca-se compreender se a
magnitude dessa degradação pode ser explicada ou prevista pelas alterações na
distribuição de energia espectral introduzidas por cada técnica.

\begin{figure}[htbp]
\centering
\resizebox{\columnwidth}{!}{%
% Requer pacotes no preâmbulo do documento principal:
% \usepackage{tikz}
% \usetikzlibrary{shapes.geometric, arrows.meta, positioning}

\begin{tikzpicture}[
    node distance=2cm and 1.5cm,
    block/.style={rectangle, draw, fill=blue!5, text width=3.5cm, text centered, rounded corners, minimum height=1.5cm, font=\sffamily\small},
    arrow/.style={-{Stealth[scale=1.2]}, thick, draw=black!70},
    dashed_arrow/.style={-{Stealth[scale=1.2]}, dashed, thick, draw=red!70!black},
    green_arrow/.style={-{Stealth[scale=1.2]}, thick, draw=green!60!black}
]

% Nós principais
\node (input) [block, fill=gray!10] {Imagem Original\\(Câmera/Sensor)};
\node (enhancement) [block, right=2.5cm of input, fill=blue!10] {Técnicas de Aprimoramento\\(Denoising, Sharpening, etc.)};
\node (output) [block, right=2.5cm of enhancement, fill=gray!10] {Imagem Aprimorada};

% Nós de desfecho
\node (human) [block, above=1.5cm of output, fill=green!10, text width=3cm] {Observador Humano\\(Qualidade Perceptual)};
\node (cnn) [block, below=1.5cm of output, fill=red!10, text width=3cm] {Classificador CNN\\(Desempenho)};

% Nó explicativo
\node (spectrum) [block, below=1.2cm of enhancement, fill=yellow!15, text width=4cm] {Alteração no Espectro\\(Perda de Altas Frequências e Sinais Discriminativos)};

% Conexões principais
\draw [arrow] (input) -- (enhancement);
\draw [arrow] (enhancement) -- (output);

% Conexões para Humano
\draw [green_arrow] (output) -- (human) node[midway, right, font=\sffamily\scriptsize, text=green!50!black] {Efeito Positivo};

% Conexões para CNN
\draw [dashed_arrow] (output) -- (cnn) node[midway, right, font=\sffamily\scriptsize, text=red!50!black] {Degradação Sistemática};

% Conexões da alteração de espectro
\draw [arrow] (enhancement) -- (spectrum) node[midway, right, font=\sffamily\scriptsize] {Causa};
\draw [dashed_arrow] (spectrum) -- (cnn) node[midway, below, font=\sffamily\scriptsize, text=red!50!black, sloped, align=center] {Redução de Features\\Relevantes};

\end{tikzpicture}
%
}
\caption{Diagrama conceitual da hipótese central: o compromisso entre a qualidade perceptual humana e a preservação de sinais discriminativos de alta frequência fundamentais para os classificadores convolucionais.}
\label{fig:hipotese_conceitual}
\end{figure}

A Seção~\ref{sec:revisao} apresenta a revisão sistemática da literatura que
fundamenta essa formulação, organizada em protocolo
(Seções~\ref{sec:qp}--\ref{sec:fb}), relatório da revisão
(Seção~\ref{sec:rr}) e respostas às questões de pesquisa
(Seção~\ref{sec:rq}).

%=======================================================================
\section{Revisão da Literatura}
\label{sec:revisao}
%=======================================================================

Esta seção reporta uma revisão sistemática da literatura conduzida para
fundamentar a pergunta central deste trabalho. O protocolo é apresentado nas
subseções \ref{sec:qp} a \ref{sec:fb}, a síntese crítica dos estudos
selecionados na subseção \ref{sec:rr}, e as respostas consolidadas às questões
de pesquisa na subseção \ref{sec:rq}.

\subsection{Questões de Pesquisa}
\label{sec:qp}

O escopo da revisão foi delimitado pelo protocolo PICOC, sumarizado na
Tabela~\ref{tab:picoc}. O objetivo é identificar, analisar e sintetizar
estudos empíricos e teóricos que investiguem o efeito de técnicas clássicas de
pré-processamento de imagem sobre a acurácia de classificadores baseados em
redes neurais convolucionais. Dá-se atenção particular às explicações
mecanísticas relacionadas à remoção de componentes de alta frequência no
domínio espectral.

\begin{table}[htbp]
\caption{Delimitação do escopo da revisão (PICOC)}
\begin{center}
\begin{tabular}{|p{1.4cm}|p{6.1cm}|}
\hline
\textbf{Elemento} & \textbf{Definição} \\
\hline
População &
Classificadores de imagem baseados em CNNs pré-treinadas, utilizadas em modo
de inferência, sem \textit{fine-tuning} ou retreinamento. \\
\hline
Intervenção &
Técnicas de pré-processamento voltadas à melhoria perceptual:
\textit{denoising}, suavização, equalização de histograma, aguçamento,
super-resolução e compressão JPEG. \\
\hline
Comparação &
Acurácia de classificação em imagens originais \textit{versus} após a
aplicação da técnica; secundariamente, \textit{benchmarks} de corrupção
padronizados. \\
\hline
Desfecho &
Acurácia de classificação, magnitude da degradação, energia espectral removida
por banda e correlação entre energia removida e queda de acurácia. \\
\hline
Contexto &
\textit{Pipelines} de pré-processamento aplicados antes da etapa de
inferência, em tarefas de classificação de imagens com CNNs. \\
\hline
\end{tabular}
\label{tab:picoc}
\end{center}
\end{table}

A partir dessa delimitação, foram formuladas quatro questões de pesquisa:

\begin{itemize}
\item \textbf{QP1.} Técnicas clássicas de pré-processamento de imagem,
aplicadas com objetivo de melhoria perceptual, degradam a acurácia de
classificadores CNN pré-treinados?
\item \textbf{QP2.} Quais mecanismos --- remoção de componentes de alta
frequência, perda de textura, deslocamento de domínio --- explicam essa
degradação?
\item \textbf{QP3.} Existe relação quantificável entre a energia espectral
removida por banda de frequência e a magnitude da queda de acurácia?
\item \textbf{QP4.} Em quais condições ou domínios de aplicação o
pré-processamento clássico melhora, em vez de degradar, o desempenho de
classificadores?
\end{itemize}

As questões são encadeadas: QP1 verifica a existência do fenômeno, QP2 busca
sua causa, QP3 testa se essa causa é quantitativamente preditiva, e QP4
delimita as fronteiras de validade do achado.

\subsection{String de Busca}
\label{sec:sb}

A \textit{string} de busca foi construída a partir dos elementos População,
Intervenção e Desfecho do PICOC, organizados em três blocos conectados por
\texttt{AND}, com sinônimos e variações terminológicas dentro de cada bloco
conectados por \texttt{OR}. O primeiro bloco delimita a população de
classificadores; o segundo, o conjunto de técnicas de pré-processamento
consideradas; o terceiro, os desfechos de interesse, incluindo tanto métricas
de desempenho quanto termos ligados à análise espectral. A \textit{string}
resultante é:

\begin{quote}
\footnotesize
\texttt{("convolutional neural network" OR "CNN" OR "DNN" OR "deep learning
classifier" OR "deep neural network")}

\texttt{AND}

\texttt{("denoising" OR "image denoising" OR "noise removal" OR "histogram
equalization" OR "CLAHE" OR "contrast enhancement" OR "image degradation" OR
"image distortion" OR "image quality degradation" OR "image preprocessing" OR
"image enhancement" OR "preprocessing techniques" OR "image sharpening" OR
"edge enhancement" OR "unsharp masking" OR "image smoothing" OR "blurring" OR
"filtering" OR "gaussian filtering" OR "low-pass" OR "super-resolution" OR
"image super resolution" OR "upsampling")}

\texttt{AND}

\texttt{("high-frequency components" OR "frequency content" OR "high frequency
information" OR "spectral energy" OR "image classification" OR "classification
accuracy" OR "classification performance" OR "robustness" OR "generalization"
OR "model robustness")}
\end{quote}

% TODO: CONFERIR a data de execucao da busca e o escopo dos campos.
% Substitua "setembro de 2026" pela data real e ajuste o escopo se a busca
% tiver sido feita em texto completo em alguma das bases.
A \textit{string} foi adaptada sintaticamente a cada fonte, preservando a
estrutura lógica dos três blocos. As buscas foram executadas em setembro de
2026 e aplicadas aos campos de título, resumo e palavras-chave, exceto no
Google Scholar, cuja interface não permite restrição de campo e indexa o texto
completo. Não foi aplicado filtro automático de ano na consulta: o critério
temporal CI2 foi verificado manualmente na etapa de triagem.

\subsection{Critérios de Inclusão}
\label{sec:ci}

Foram incluídos na revisão os estudos que satisfizessem simultaneamente:

\begin{itemize}
\item \textbf{CI1.} Estudos disponíveis em inglês ou português, com texto
completo acessível.
\item \textbf{CI2.} Estudos publicados a partir de 2015.
\item \textbf{CI3.} Estudos publicados em conferências ou periódicos revisados
por pares.
\item \textbf{CI4.} Estudos que avaliem quantitativamente o efeito de técnicas
de pré-processamento de imagem sobre a acurácia de classificadores baseados em
redes neurais.
\end{itemize}

Foi admitida uma exceção controlada ao critério CI3: pré-prints de alta
relevância metodológica e amplamente citados na literatura da área foram
incluídos quando forneciam evidência quantitativa com \textit{baseline}
explícito, com a ressalva de que o veículo de publicação revisado por pares não
é informado. É o caso, entre os estudos selecionados, de
\cite{b4}, \cite{b5}, \cite{b8}, \cite{b9}, \cite{b10}, \cite{b13},
\cite{b14}, \cite{b22}, \cite{b24}, \cite{b26} e \cite{b27}. O estudo
\cite{b24} constitui o caso mais extremo dessa exceção, por tratar-se de
manuscrito não publicado e sem indicação de veículo; sendo utilizado apenas como
evidência convergente, nunca como sustentação isolada de qualquer afirmação.

Dois estudos foram admitidos em estatuto secundário por avaliarem detecção de
objetos ou segmentação, e não classificação direta: \cite{b13} e parte dos resultados
de \cite{b12}. Toda afirmação derivada deles é explicitamente rotulada como tal
ao longo da síntese. Um estudo foi admitido como referência teórica, sem contar
como evidência empírica para nenhuma questão: \cite{b3}.

\subsection{Critérios de Exclusão}
\label{sec:ce}

Foram excluídos os estudos que atendessem a qualquer um dos critérios:

\begin{itemize}
\item \textbf{CE1.} Duplicatas ou versões preliminares de estudos já incluídos.
\item \textbf{CE2.} Estudos focados exclusivamente em tarefas que não sejam
classificação de imagem --- por exemplo, apenas detecção de objetos ou
segmentação, sem métrica de classificação.
\item \textbf{CE3.} Estudos que tratam exclusivamente de ataques adversariais
intencionais, fora do escopo de melhoria perceptual.
\item \textbf{CE4.} Estudos sem avaliação quantitativa de acurácia ou
desempenho de classificação.
\item \textbf{CE5.} Trabalhos não revisados por pares --- pré-prints não
publicados, \textit{blogs}, relatórios técnicos --- exceto quando citados
apenas como motivação e não como evidência formal.
\end{itemize}

A Tabela~\ref{tab:funil} apresenta o funil de seleção completo, e a
Tabela~\ref{tab:excluidos} registra os estudos excluídos na etapa de triagem
por leitura integral, com o critério acionado em cada caso. O registro é
mantido porque três deles contêm fundamentação conceitual relevante e poderiam
ser reconsiderados caso o escopo da QP2 fosse ampliado.

% TODO (1 de 2 pendencias): PREENCHER as QUATRO linhas marcadas com "--".
% Os numeros 76, 5, 29, 27, 1 e 1 ja estao conferidos; faltam apenas os
% intermediarios (duplicatas removidas, submetidos e excluidos na triagem por
% titulo/resumo, submetidos a leitura integral). Nao estao na planilha
% data_extraction.xls -- pegue no relatorio da ferramenta de apoio.
\begin{table}[htbp]
\caption{Funil de seleção dos estudos}
\begin{center}
\footnotesize
\begin{tabular}{|p{5.3cm}|c|}
\hline
\textbf{Etapa} & \textbf{Estudos} \\
\hline
Importados das fontes de busca & 76 \\
\hline
Removidos por duplicação (CE1) & -- \\
\hline
Submetidos à triagem por título e resumo & -- \\
\hline
Excluídos na triagem por título e resumo & -- \\
\hline
Submetidos à leitura integral & -- \\
\hline
Excluídos na leitura integral & 5 \\
\hline
\textbf{Incluídos na revisão} & \textbf{29} \\
\hline
\quad evidência empírica primária & 27 \\
\hline
\quad evidência secundária (detecção) & 1 \\
\hline
\quad referência teórica & 1 \\
\hline
\end{tabular}
\label{tab:funil}
\end{center}
\end{table}

\begin{table}[htbp]
\caption{Estudos excluídos na triagem por leitura integral}
\begin{center}
\footnotesize
\begin{tabular}{|p{3.3cm}|p{0.7cm}|p{2.9cm}|}
\hline
\textbf{Estudo} & \textbf{Crit.} & \textbf{Justificativa} \\
\hline
Nakkiran (2019), \textit{Adversarial robustness may be at odds with
simplicity} & CE4 & Trabalho inteiramente teórico sobre distribuições
sintéticas; sem imagens, CNN, \textit{dataset} ou experimento. \\
\hline
Jacobsen \textit{et al.} (2019), \textit{Excessive invariance causes
adversarial vulnerability} & CE3 & Estudo centrado em vulnerabilidade
adversária por invariância excessiva; o pré-processamento não é variável do
estudo e não há análise espectral. \\
\hline
Luo \textit{et al.} (2023), \textit{Spectral discriminators for perceptual
quality improvement} & CE4 & Desfechos são PSNR, SSIM, LPIPS e correlação com
juízo humano; sem acurácia. \\
\hline
Sheng \textit{et al.} (2023), \textit{Frequency-domain deep guided image
denoising} & CE4 & Estudo exclusivamente de restauração; ``\textit{accuracy}''
refere-se à fidelidade da reconstrução. \\
\hline
Ayyoubzadeh e Wu (2021), \textit{High frequency detail accentuation in CNN
image restoration} & CE4 & Restauração avaliada por PSNR, SSIM e NIQE; nenhum
classificador é avaliado. \\
\hline
\end{tabular}
\label{tab:excluidos}
\end{center}
\end{table}

\subsection{Fontes de Busca}
\label{sec:fb}

A busca foi executada sobre sete bases digitais, escolhidas por cobrirem de
forma complementar as comunidades de visão computacional, aprendizado de
máquina e processamento de sinais: ACM Digital Library
(\texttt{portal.acm.org}), arXiv (\texttt{arxiv.org}), Google Scholar
(\texttt{scholar.google.com}), IEEE Xplore Digital Library
(\texttt{ieeexplore.ieee.org}), ScienceDirect (\texttt{sciencedirect.com}),
Scopus (\texttt{scopus.com}) e SpringerLink (\texttt{link.springer.com}).

A Tabela~\ref{tab:fontes} apresenta a quantidade de estudos importados por
fonte. A concentração no arXiv e no IEEE Xplore reflete uma característica do
campo: a produção sobre robustez de CNNs a degradações circula
predominantemente em pré-prints e em anais de conferências de primeira linha
(CVPR, ICLR, NeurIPS, ICML, ICASSP), cuja indexação nas demais bases é parcial
ou tardia. As buscas em ACM, ScienceDirect, Scopus e SpringerLink não
retornaram resultados adicionais que satisfizessem os critérios de inclusão e
que já não tivessem sido recuperados pelas duas primeiras fontes.

\begin{table}[htbp]
\caption{Estudos importados por fonte de busca}
\begin{center}
\begin{tabular}{|l|c|}
\hline
\textbf{Fonte} & \textbf{Estudos importados} \\
\hline
arXiv & 36 \\
IEEE Xplore Digital Library & 38 \\
Google Scholar & 2 \\
ACM Digital Library & 0 \\
ScienceDirect & 0 \\
Scopus & 0 \\
SpringerLink & 0 \\
\hline
\textbf{Total} & \textbf{76} \\
\hline
\end{tabular}
\label{tab:fontes}
\end{center}
\end{table}

Dos 76 estudos importados, 29 foram selecionados após a aplicação dos critérios
de inclusão e exclusão. Destes, 28 fornecem evidência empírica --- um deles
\cite{b13} apenas como evidência secundária de detecção de objetos --- e um
\cite{b3} foi admitido como referência teórica. A
Tabela~\ref{tab:caracterizacao} caracteriza o corpus resultante. A extração
formal de dados, conduzida em formulário estruturado, cobre 20 desses 29
estudos; os nove restantes foram incorporados à síntese a partir da leitura
integral, sem preenchimento do formulário, e estão identificados na
Tabela~\ref{tab:caracterizacao}.

Os estudos submetidos à extração formal foram avaliados por um instrumento de
qualidade de seis questões, respondidas como \textit{Sim},
\textit{Parcialmente} ou \textit{Não}, cobrindo: reporte de métricas
quantitativas, discussão de mecanismos explicativos, realização de análise
no domínio da frequência, uso de múltiplos \textit{datasets}, uso de múltiplas
arquiteturas e disponibilização de código ou dados. A Tabela~\ref{tab:qualidade}
consolida as respostas registradas, usadas ao longo da síntese para ponderar o peso
atribuído a cada evidência, e a Tabela~\ref{tab:extracao} reporta a
distribuição dos demais campos do formulário.

\begin{table}[htbp]
\caption{Instrumento de qualidade: distribuição das respostas registradas ($n=20$)}
\begin{center}
\footnotesize
\begin{tabular}{|p{4.5cm}|c|c|c|}
\hline
\textbf{Critério de qualidade} & \textbf{Sim} & \textbf{Parc.} & \textbf{Não} \\
\hline
Q1. Métricas quantitativas reportadas$^{\ast}$ & 20 & 0 & 0 \\
\hline
Q2. Discussão de mecanismos explicativos$^{\ast}$ & 20 & 0 & 0 \\
\hline
Q3. Análise no domínio da frequência & 12 & 2 & 6 \\
\hline
Q4. Múltiplos \textit{datasets}$^{\dagger}$ & 16 & 0 & 4 \\
\hline
Q5. Múltiplas arquiteturas & 16 & 0 & 4 \\
\hline
Q6. Código ou dados disponíveis & 10 & 0 & 10 \\
\hline
\multicolumn{4}{|p{6.9cm}|}{\footnotesize $^{\ast}$ Derivado da presença de valor no campo correspondente, e não de julgamento graduado: todos os 20 registros informam ao menos uma métrica (acurácia em 18, F1-score em 1, outras em 13) e descrevem um mecanismo explicativo. $^{\dagger}$ Derivado da contagem de \textit{datasets} declarados por estudo.} \\
\hline
\end{tabular}
\label{tab:qualidade}
\end{center}
\end{table}

\begin{table}[htbp]
\caption{Distribuição dos demais campos da extração de dados ($n=20$)}
\begin{center}
\footnotesize
\begin{tabular}{|p{3.6cm}|p{3.3cm}|}
\hline
\textbf{Campo} & \textbf{Distribuição} \\
\hline
Direção do efeito sobre a acurácia & Efeito misto 14; piora 5; melhora 1 \\
\hline
Intensidade avaliada em múltiplos níveis & Sim 17; não 3 \\
\hline
Métrica de desempenho & Acurácia 18; outras 13; F1-score 1 \\
\hline
Ano de publicação & 2016--2018: 6; 2019--2021: 10; 2022--2026: 4 \\
\hline
\end{tabular}
\label{tab:extracao}
\end{center}
\end{table}

% TODO (2 de 2 pendencias): restam cinco celulas "n.i.", nas linhas de b5, b9,
% b26 e b27 -- os quatro estudos que nao constam da planilha de extracao.
% Todas as demais foram preenchidas com data_extraction.xls.
\begin{table*}[htbp]
\caption{Caracterização dos estudos incluídos. Esp.: realiza análise no domínio da frequência. Efeito: direção do efeito do pré-processamento sobre o desempenho reportado.}
\begin{center}
\footnotesize
\setlength{\tabcolsep}{3pt}
\begin{tabular}{|c|p{3.1cm}|p{2.7cm}|p{3.1cm}|p{1.7cm}|c|p{1.4cm}|p{1.3cm}|}
\hline
\textbf{Ref.} & \textbf{Técnica(s) avaliada(s)} & \textbf{Arquitetura(s)} &
\textbf{Dataset(s)} & \textbf{Métrica} & \textbf{Esp.} & \textbf{Efeito} &
\textbf{QPs} \\
\hline
\multicolumn{8}{|l|}{\footnotesize Linhas com \emph{n.i.} e as de
\cite{b1}, \cite{b2}, \cite{b3}, \cite{b5}, \cite{b9}, \cite{b11},
\cite{b13}, \cite{b26} e \cite{b27} não passaram pelo formulário de extração;
seus valores derivam da leitura integral.} \\
\hline
\cite{b1} & \textit{Blur}, ruído, compressão, contraste & Quatro CNNs pré-treinadas &
ImageNet (validação) & Acurácia & Não & Piora & QP1, QP2 \\
\hline
\cite{b2} & \textit{Blur}, ruído, oclusão, escurecimento & Três CNNs pré-treinadas &
LFW & Acurácia \textit{rank}-1 & Não & Piora & QP1, QP2 \\
\hline
\cite{b3} & Super-resolução, \textit{denoising}, \textit{deblurring} &
n.a.\ (resultado teórico) & BSD100 & PSNR, SSIM, NIQE & Não & n.a. &
Teórica \\
\hline
\cite{b4} & Desfoque, \textit{blur} gaussiano, \textit{blur} de movimento & VGG-16 & ImageNet &
Acurácia \textit{top}-5 & Sim & Piora & QP1--QP4 \\
\hline
\cite{b5} & Seis modelos de degradação & Seis arquiteturas & n.i. & Acurácia &
Não & Piora & QP1 \\
\hline
\cite{b6} & \textit{Non-Local Means}, CLAHE & ResNet-50 & ImageNet, ImageNet-C, ImageNet-P, CIFAR-10-C
& Erro, mCE & Não & Misto & QP1--QP4 \\
\hline
\cite{b7} & Autoencoder \textit{denoising} & Inception-v3 & ImageNet &
Acurácia, outras & Não & Misto & QP1--QP4 \\
\hline
\cite{b8} & Filtro de mediana, \textit{Non-Local Means} & Não especificado
(SVM sobre LBP e HOG) & Corel, Caltech101-600 & F1-score & Não & Misto &
QP1, QP2, QP4 \\
\hline
\cite{b9} & Filtragem como pré-processamento e como canal adicional &
Inception-v3 & n.i. & Acurácia & Não & Misto & QP1, QP4 \\
\hline
\cite{b10} & Compressão JPEG & Outra (não especificada) & ImageNet
(ILSVRC 2012) & Acurácia, outras & Não & Misto & QP1, QP2, QP4 \\
\hline
\cite{b11} & Deflexão de \textit{pixels} e limiar \textit{wavelet} &
ResNet-50 & ImageNet & Acurácia & Sim & Piora leve & QP1, QP4 \\
\hline
\cite{b12} & Reconstrução por ISP (RAW$\rightarrow$RGB) & ResNet-18,
CenterNet, QueryInst & CIFAR-100, Zurich, LOD, LOD-Seg, SIDD & Acurácia, mAP & Sim &
Misto & QP1, QP2, QP4 \\
\hline
\cite{b13} & Resolução espacial, compressão BPG & RA-YOLO (YOLOv5) & TJU-DHD,
EuroCity, BDD100K & mAP & Não & Misto & QP1, QP3 \\
\hline
\cite{b14} & Máscaras radiais e aleatórias de Fourier &
ResNet & SVHN, CIFAR-10 & Erro, outras & Sim & Piora & QP1--QP4 \\
\hline
\cite{b15} & Filtragem LFC/HFC, suavização de \textit{kernels} & ResNet-18,
VGG, AlexNet, LeNet & CIFAR-10, MNIST, FashionMNIST, ImageNet, COCO & Acurácia, mAP
& Sim & Misto & QP1, QP2, QP3 \\
\hline
\cite{b16} & Conflito textura/forma, Stylized-ImageNet & CNNs ImageNet e
observadores humanos & ImageNet, Stylized-ImageNet & Acurácia \textit{top}-1 &
Parcial & Misto & QP1--QP4 \\
\hline
\cite{b17} & Remoção de características não robustas & ResNet, VGG, outra &
CIFAR-10, ImageNet (subconjunto) & Acurácia & Não & Misto & QP1, QP2, QP4 \\
\hline
\cite{b18} & Corrupções do \textit{benchmark} e augmentação & ResNet, outra &
CIFAR-10/-C, ImageNet/-C & Acurácia, outras & Sim & Misto & QP1--QP4 \\
\hline
\cite{b19} & Perturbação por base de Fourier & MLP, LeNet, WideResNet,
DenseNet, VGG, ResNet-50, GoogLeNet & MNIST, FashionMNIST, SVHN, CIFAR-10,
CIFAR-100, ILSVRC2015 & \textit{Fool ratio} & Sim & Piora & QP2, QP3 \\
\hline
\cite{b20} & Regularização de viés de frequência & ResNet, outra & ImageNet-C,
ImageNet-R, ObjectNet & Erro, outras & Sim & Misto & QP1--QP4 \\
\hline
\cite{b21} & \textit{Blur} gaussiano, JPEG, ruído, recorte & CNN rasa, \textit{ridge},
kNN, Eigenfaces & FFHQ, LSUN, CelebA, Stanford Dogs & Acurácia & Sim & Piora &
QP1--QP4 \\
\hline
\cite{b22} & Mascaramento probabilístico de bandas DCT & ResNet-18, VGG-16 &
CIFAR-10, TinyImageNet, ImageNet & Acurácia & Sim & Misto & QP1--QP4 \\
\hline
\cite{b23} & Filtragem passa-baixa em treino e em teste & ResNet, outra &
MNIST, CIFAR-10, ImageNet, sintéticos & Acurácia, outras & Sim & Misto &
QP1--QP4 \\
\hline
\cite{b24} & Filtragem espectral, embaralhamento de fase & ResNet-18 &
ImageNet-50 & Acurácia, outras & Sim & Misto & QP1, QP2, QP4 \\
\hline
\cite{b25} & \textit{Blur} gaussiano, ruído branco & AlexNet, ResNet & ImageNet,
Caltech-101, Caltech-256, SUN-397 & Acurácia, outras & Não & Piora &
QP1--QP4 \\
\hline
\cite{b26} & Autoencoder \textit{denoising} & VGG-16, Inception-v3 & n.i. &
Acurácia & Não & Misto & QP2, QP4 \\
\hline
\cite{b27} & Aproximação de baixo posto & n.i. & n.i. & Acurácia & Sim &
Misto & QP2 \\
\hline
\cite{b28} & \textit{Blur} gaussiano, ruído, blocagem, \textit{Frequency Dropout} &
\textit{Tiny} ResNet & CIFAR-10, CIFAR-100 & Acurácia & Sim & Misto & QP1,
QP4 \\
\hline
\cite{b29} & \textit{Random Erasing} no domínio da frequência & ResNet-20 a
ResNet-110 & CIFAR-10/-C, CIFAR-100/-C & Acurácia & Parcial & Melhora & QP4 \\
\hline
\end{tabular}
\label{tab:caracterizacao}
\end{center}
\end{table*}

A Tabela~\ref{tab:cobertura} resume a cobertura das técnicas de interesse,
tabulada sobre o vocabulário controlado do formulário de extração. Ela torna
explícita uma assimetria relevante para o posicionamento deste trabalho:
suavização e \textit{denoising} concentram a maior parte da evidência, a
equalização de histograma é registrada em um único estudo, e aguçamento e
super-resolução não são registrados em nenhum --- apesar de constarem como
opções do formulário.

\begin{table}[htbp]
\caption{Cobertura das técnicas de interesse, conforme o vocabulário controlado do formulário de extração ($n=20$; um estudo pode registrar mais de uma técnica)}
\begin{center}
\footnotesize
\begin{tabular}{|p{3.4cm}|c|c|}
\hline
\textbf{Técnica} & \textbf{$n$} & \textbf{\% dos 18$^{\ast}$} \\
\hline
Outra (não catalogada) & 14 & 78\% \\
\hline
Suavização / \textit{blur} & 11 & 61\% \\
\hline
Compressão JPEG & 5 & 28\% \\
\hline
\textit{Denoising} & 5 & 28\% \\
\hline
Equalização de histograma & 1 & 6\% \\
\hline
Aguçamento / \textit{sharpening} & 0 & 0\% \\
\hline
Super-resolução & 0 & 0\% \\
\hline
\multicolumn{3}{|p{6.9cm}|}{\footnotesize $^{\ast}$ Dois dos 20 registros estão
com o campo de técnica em branco. As sete categorias listadas são exatamente as
do vocabulário controlado do formulário, de modo que a contagem nula de
aguçamento e de super-resolução é um resultado da extração, e não uma omissão.} \\
\hline
\end{tabular}
\label{tab:cobertura}
\end{center}
\end{table}

\subsection{Relatório da Revisão}
\label{sec:rr}

\subsubsection{Qualidade perceptual e qualidade para a máquina}

A suposição implícita em boa parte dos \textit{pipelines} de visão
computacional é a de que uma imagem melhor para o observador humano é também
uma imagem melhor para o classificador. A literatura revisada mostra que essa
equivalência não se sustenta. Dodge e Karam \cite{b1} foram os primeiros a
formular o problema em termos explícitos, distinguindo a qualidade perceptual
humana da qualidade do ponto de vista de um sistema automático. Eles argumentaram
que distorções cotidianas de aquisição, transmissão e armazenamento são muito
mais prováveis na prática do que perturbações adversárias construídas deliberadamente. 

Avaliando quatro arquiteturas pré-treinadas sobre o conjunto de validação do ImageNet 
degradado por cinco famílias de distorção, os autores documentaram uma assimetria 
que reaparece de forma sistemática em toda a literatura posterior: as redes são 
muito sensíveis a borramento e ruído, mas notavelmente resilientes a compressão moderada 
e alterações de contraste. Mais relevante para o presente trabalho, registram que, 
mesmo em desvios-padrão de ruído nos quais o desempenho médio das redes cai abaixo de 20\%, 
as imagens permanecem facilmente reconhecíveis por observadores humanos.

Karahan \textit{et al.} \cite{b2} replicam essa dissociação em reconhecimento
facial e a tornam ainda mais aguda. Aplicando cinco famílias de degradação
apenas ao conjunto de sondas do protocolo de identificação fechada do LFW,
observam que a acurácia \textit{rank}-1 de três CNNs pré-treinadas cai a menos
de 6\% sob níveis de ruído nos quais o olho humano ainda identifica os 
sujeitos sem dificuldade. No sentido inverso, verificam que as redes são surpreendentemente 
resilientes ao escurecimento, uma distorção que prejudica o observador humano. As duas escalas 
de qualidade, portanto, não apenas divergem em magnitude, mas também divergem em ordenação.

Há também uma razão formal para esperar tensão entre aparência e fidelidade.
Blau e Michaeli \cite{b3} demonstram que qualidade perceptual e distorção são
objetivos mutuamente incompatíveis, apresentando uma fronteira de compromisso
monotonicamente não crescente e convexa. Esse resultado é válido para qualquer
medida de distorção --- inclusive distâncias definidas em espaços de
características de redes neurais. Convém registrar, contudo, que esse trabalho
não mede acurácia de classificação: o ``classificador'' do teorema é um
discriminador ótimo entre imagens reais e sintéticas, e a parte experimental
compara algoritmos de super-resolução por métricas de qualidade de imagem. A
ponte entre o compromisso percepção--distorção e o desempenho de reconhecimento
é, portanto, uma inferência que a presente revisão assume explicitamente como
hipótese, e não um resultado estabelecido. Sua utilidade é delimitar a
expectativa teórica: se otimizar a aparência necessariamente afasta a imagem
reconstruída da imagem original, e se o classificador depende de propriedades
dessa imagem original, então a degradação de desempenho tem uma explicação
anterior a qualquer análise espectral.

\subsubsection{Evidência empírica da degradação}

O corpo empírico revisado converge para a conclusão de que as transformações
clássicas de entrada cobram um preço em acurácia, mesmo quando aplicadas a
imagens que não sofreram qualquer degradação prévia. A magnitude desse preço,
porém, varia por ordens de grandeza conforme a técnica, o \textit{dataset} e o
protocolo. A Tabela~\ref{tab:magnitudes} consolida as magnitudes
reportadas com \textit{baseline} explícito.

\begin{table}[htbp]
\caption{Efeito reportado sobre imagens não degradadas previamente, em estudos com \textit{baseline} explícito}
\begin{center}
\footnotesize
\begin{tabular}{|p{2.0cm}|p{2.0cm}|p{2.6cm}|}
\hline
\textbf{Técnica} & \textbf{Modelo / dados} & \textbf{Efeito reportado} \\
\hline
Desfoque, raio 8 \cite{b4} & VGG-16 / ImageNet & \textit{top}-5: 90,88\%
$\rightarrow$ 60,97\% \\
\hline
\textit{Blur} gaussiano, $\sigma=8$ \cite{b4} & VGG-16 / ImageNet & \textit{top}-5:
90,88\% $\rightarrow$ 17,39\% \\
\hline
\textit{Blur} gaussiano $3\times3$ \cite{b28} & \textit{Tiny} ResNet / CIFAR-10 &
84,72\% $\rightarrow$ 67,84\% \\
\hline
Ruído gaussiano \cite{b28} & \textit{Tiny} ResNet / CIFAR-10 & 84,72\%
$\rightarrow$ 61,35\% \\
\hline
Blocagem por reamostragem \cite{b28} & \textit{Tiny} ResNet / CIFAR-10 &
84,72\% $\rightarrow$ 64,28\% \\
\hline
Corrupções comuns \cite{b29} & ResNet-44 / CIFAR-10-C & erro: 5,53\%
$\rightarrow$ 28,35\% \\
\hline
Autoencoder \textit{denoising} \cite{b7} & Inception-v3 / ImageNet & 76,7\%
$\rightarrow$ 58,3\% \\
\hline
\textit{Non-Local Means} \cite{b6} & ResNet-50 / ImageNet-C & mCE: 76,7\%
$\rightarrow$ 82,1\% \\
\hline
Passa-baixa só no teste \cite{b23} & CIFAR-10 & $\approx$93\% $\rightarrow$
27,45\% \\
\hline
Máscara radial de Fourier \cite{b14} & ResNet-200 / CIFAR-10 & erro: 5,22\%
$\rightarrow$ 23,37\% \\
\hline
\textit{Anti-aliasing} crescente \cite{b21} & Ridge-DCT / LSUN & 98,24\%
$\rightarrow$ 84,20\% \\
\hline
Deflexão $+$ \textit{wavelet} \cite{b11}$^{\dagger}$ & ResNet-50 / ImageNet &
100\% $\rightarrow$ 98,3\% \\
\hline
CLAHE \cite{b6} & ResNet-50 / ImageNet & erro: 23,87\% $\rightarrow$ 23,55\%
(melhora) \\
\hline
\multicolumn{3}{|p{6.9cm}|}{\footnotesize $^{\dagger}$ O \textit{baseline} de
100\% decorre do protocolo do estudo, que mede o efeito da transformação apenas
sobre o subconjunto de imagens originalmente classificadas de forma correta.} \\
\hline
\end{tabular}
\label{tab:magnitudes}
\end{center}
\end{table}

O borramento é o caso mais bem documentado. Vasiljevic \textit{et al.}
\cite{b4} avaliam a VGG-16 pré-treinada sobre versões borradas do conjunto de
validação do ImageNet e reportam queda monotônica da acurácia \textit{top}-5 de
90,88\% para 88,06\%, 81,48\%, 72,69\% e 60,97\% com desfoque de raio 2, 4, 6 e
8. Um detalhe importante para a presente investigação é que o
borramento de movimento de comprimento comparável é substancialmente menos
danoso (83,32\% contra 60,97\%). Os autores atribuem essa diferença ao fato de a
atenuação de alta frequência ser unidirecional em um caso e radialmente
uniforme no outro --- a mesma ``quantidade'' de suavização produz danos
distintos conforme a geometria da supressão espectral. Roy \textit{et al.}
\cite{b5} ampliam a varredura para seis arquiteturas contra seis modelos de
degradação e encontram decaimento monotônico em todos os casos, notando que 
as arquiteturas mais rasas resistem melhor do que as mais profundas. Wu \cite{b28} 
replica a mesma ordem de magnitude em escala menor e com \textit{baseline} limpo 
explícito, reportando queda de 84,72\% para 67,84\% sob \textit{blur} gaussiano, 
64,28\% sob blocagem e 61,35\% sob ruído gaussiano no CIFAR-10.

O \textit{denoising} --- a técnica cuja finalidade declarada é justamente
melhorar a imagem --- produz os resultados mais contraintuitivos. Hendrycks e
Dietterich \cite{b6} aplicam \textit{Non-Local Means} como etapa de
pré-processamento sobre o ImageNet-C e registram que o erro médio de corrupção
sobe de 76,7\% para 82,1\%: a restauração torna o classificador globalmente
pior. Liao \textit{et al.} \cite{b7} quantificam o custo diretamente sobre
imagens não atacadas, mostrando que um autoencoder \textit{denoising} padrão
derruba a acurácia do Inception-v3 de 76,7\% para 58,3\% --- 18,4 pontos
percentuais perdidos sem que houvesse qualquer degradação a remover. Paranhos
da Costa \textit{et al.} \cite{b8} observam que filtragem de mediana e 
\textit{Non-Local Means} recuperam parte do desempenho perdido sobre entradas 
ruidosas, mas jamais alcançam o \textit{baseline} limpo. Eles atribuem explicitamente 
a lacuna residual à perda de detalhe e textura causada pela suavização. Yim e Sohn \cite{b9} 
chegam a uma conclusão negativa equivalente com uma Inception-v3: o pré-processamento 
isolado não produz melhoria consistente e o modelo alimentado apenas com a imagem 
filtrada perde desempenho.

A compressão e a quantização exibem um padrão qualitativamente diferente, de
platô seguido de queda. Dziugaite \textit{et al.} \cite{b10} fornecem o dado
primário de que a compressão JPEG reduz a confiança média do classificador 
mesmo sobre imagens limpas. Prakash \textit{et al.} \cite{b11} elevam essa observação 
a diagnóstico do campo, afirmando que o maior obstáculo às defesas baseadas em 
transformação da entrada é precisamente o fato de a transformação degradar a qualidade 
das imagens não adversárias. Ainda assim, ambos os trabalhos e também \cite{b1} registram que a
compressão só se torna danosa abaixo de um limiar de qualidade relativamente
baixo --- um comportamento não linear que contrasta com a degradação monotônica
observada sob borramento.

Por fim, dois trabalhos recentes deslocam a evidência para \textit{pipelines}
de reconstrução de estado da arte. Chen e Ma \cite{b12} documentam que
algoritmos de ISP inteligente projetados para melhorar a qualidade da imagem
podem produzir desempenho de reconhecimento inferior ao do ISP tradicional da
câmera. O experimento mais valioso da equipe é de controle: fixando a qualidade de
reconstrução em 21,30\,dB de PSNR, ainda obtêm 2,0 mAP de ganho sobre as
imagens RGB tradicionais --- ou seja, o desempenho varia com o conteúdo
espectral da reconstrução mesmo quando a qualidade perceptual é mantida
constante. Hao \textit{et al.} \cite{b13}, em evidência secundária de detecção
de objetos, reforçam o mesmo descolamento ao observar que é possível comprimir
imagens em amplitude de forma substancial sem afetar a acurácia de detecção,
mesmo quando a qualidade visual já se deteriorou.

\subsubsection{A hipótese espectral}

A explicação que melhor articula esse conjunto de resultados sustenta que
classificadores convolucionais exploram, como sinal discriminativo,
componentes da imagem que o sistema visual humano descarta --- e que as
técnicas de aprimoramento perceptual, por construção, atuam exatamente sobre
esses componentes.

A formulação operacional mais antiga é a de Jo e Bengio \cite{b14}, que propõem
medir se as CNNs aprendem conceitos semânticos ou apenas regularidades
estatísticas superficiais. Aplicando máscaras radiais e aleatórias no domínio
de Fourier, calibradas visualmente para preservar a reconhecibilidade humana
dos objetos, os autores observam que o erro de uma \textit{Preact}-ResNet-200
sobe de 5,22\% para 23,37\% no CIFAR-10. A conclusão é que a rede depende de
energia espectral da qual o observador humano prescinde. Wang \textit{et al.} \cite{b15} 
tornam a hipótese explícita e a ancoram em uma série de experimentos: 
enquanto humanos dependem de componentes semânticos de baixa frequência, as CNNs
minimizam a perda explorando também sinais imperceptíveis de alta frequência.
Mostram ainda que forçar a suavidade dos \textit{kernels} convolucionais 
--- o que equivale a impor um filtro passa-baixa --- derruba a acurácia em 
dados limpos de 85,6\% para 63,1\%.

Geirhos \textit{et al.} \cite{b16} fornecem a verificação causal mais forte do
mecanismo, ainda que formulada no domínio espacial. Com 48.560 ensaios
psicofísicos, estabelecem que CNNs treinadas no ImageNet classificam por 
textura local enquanto humanos classificam por forma global. Se a decisão da 
rede se apoia em textura, então qualquer transformação que suavize textura 
remove o próprio sinal de decisão --- explicação oferecida de forma independente 
por \cite{b1}, \cite{b2} e \cite{b8}. Ilyas \textit{et al.} \cite{b17} complementam 
o argumento pelo lado da utilidade: características imperceptíveis ao humano são 
genuinamente preditivas, e removê-las custa cerca de dez pontos de acurácia padrão 
no CIFAR-10. Contudo, a ponte explícita entre ``características imperceptíveis'' e 
``alta frequência'' é feita posteriormente por \cite{b18}.

A sensibilidade espectral tem ainda uma origem arquitetural. Tsuzuku e Sato
\cite{b19} derivam que os vetores singulares de camadas convolucionais 
são combinações de poucas funções de base de Fourier, verificando empiricamente 
que existem ``pontos sensíveis'' bem definidos no espectro nos quais uma 
perturbação de amplitude fixa altera a predição. O controle metodológico decisivo 
é que um MLP treinado sobre os mesmos dados não exibe esses pontos: a sensibilidade 
é atribuível à operação de convolução. 

\subsubsection{Da dependência espectral à previsão da degradação}

O passo seguinte --- estabelecer uma relação quantificável entre energia
espectral alterada e magnitude da queda de acurácia --- é justamente onde a
literatura se torna mais escassa e menos convergente.

A evidência mais próxima de uma relação dose--resposta explícita é a de Yin
\textit{et al.} \cite{b18}. Os autores caracterizam cada corrupção do
CIFAR-10-C e do ImageNet-C por uma métrica escalar (a fração de energia de alta
frequência do delta introduzido pela corrupção) e relacionam esse valor à
variação de acurácia por regressão linear. Seus mapas de calor mostram que o 
modelo treinado naturalmente é sensível a perturbações em quase todo o espectro, 
exceto nas frequências mais baixas. Eles também estabelecem uma assimetria estrutural: 
o desempenho cai muito mais rapidamente quando se remove informação de baixa
frequência do que de alta, porque imagens naturais concentram sua energia nas
baixas frequências.

Saikia \textit{et al.} \cite{b20} convertem essa ideia em um eixo ordinal,
ordenando as quinze corrupções do ImageNet-C por conteúdo de alta frequência, e
quantificam o custo dessa escolha em acurácia limpa. Frank \textit{et al.}
\cite{b21} oferecem a curva dose--resposta mais limpa em um domínio específico,
mostrando que a acurácia de um classificador cai de 98,24\% para 85,96\% e 
84,20\% à medida que o \textit{kernel} de \textit{anti-aliasing} cresce --- 
uma relação direta entre intensidade de suavização espectral e desempenho.

Contra a leitura simples dessa hipótese pesam, contudo, quatro resultados
importantes. Primeiro, Maiya \textit{et al.} \cite{b22} mostram que a banda 
crítica depende do conjunto de dados e não é universalmente a alta frequência: 
remover 90\% da banda 0--15 custa cerca de 2\% de acurácia no CIFAR-10 e 35\% no 
ImageNet. Segundo, Tsuzuku e Sato \cite{b19} encontram que as bandas mais danosas 
são as baixas e médias, não as altas. Terceiro, Ortiz-Jiménez \textit{et al.}
\cite{b23} demonstram que a magnitude da degradação depende menos da energia
removida do que de \emph{onde} ela é removida no protocolo: filtrar apenas o 
conjunto de teste derruba a acurácia do CIFAR-10 de $\approx$93\% para 27,45\%, 
enquanto treinar com os mesmos dados filtrados custa cerca de três pontos. 
Quarto, o próprio \cite{b18} reporta que treinar com ruído de estatísticas de Fourier 
idênticas às da corrupção de névoa piora o desempenho sob névoa; e Shao \textit{et al.}
\cite{b24} observam que a filtragem tradicional sequer produz o deslocamento
espectral limpo que se lhe atribui.

\subsubsection{Mecanismos concorrentes e complementares}

A explicação espectral não é a única disponível, e os trabalhos que propõem
alternativas ajudam a delimitar o que ela efetivamente explica.

O primeiro é o descasamento de distribuição. Vasiljevic \textit{et al.}
\cite{b4} mostram que o ajuste fino com imagens borradas quase fecha a lacuna
de desempenho, o que permite descartar tanto a deficiência arquitetural quanto
a ausência de informação discriminativa: o sinal semântico sobrevive, mas 
passa a ser mal representado. Paranhos da Costa \textit{et al.} \cite{b8}
chegam à mesma conclusão, observando que o pré-processamento cria um 
novo domínio em vez de restaurar o original.

Um segundo mecanismo opera no nível do filtro individual. Borkar e Karam
\cite{b25} mostram que corrigir apenas metade dos filtros em uma única 
camada já recupera parte substancial do desempenho. Os filtros mais atingidos 
pelo borramento gaussiano são precisamente os seletivos em frequência e orientação (detectores de borda e contorno).

Um terceiro mecanismo, proposto por Liao \textit{et al.} \cite{b7}, é o efeito
de amplificação de erro. O resíduo deixado por um \textit{denoiser}, 
embora pequeno na entrada, é progressivamente amplificado ao longo da 
hierarquia até atingir magnitude de perturbação adversária no topo. 
Isso explica por que menor erro de reconstrução em espaço de \textit{pixels} não implica 
melhor classificação \cite{b26}. 

Convém registrar, finalmente, que parte da literatura de defesa adversária
opera sobre a premissa inversa --- a de que a perturbação relevante é de alta
frequência e que descartá-la é, portanto, seguro \cite{b27}. É exatamente essa
premissa que os trabalhos com instrumentação espectral direta \cite{b18},
\cite{b22} contestam.

\subsubsection{Condições em que o pré-processamento melhora}

A fronteira mais nítida é traçada por Hendrycks e Dietterich \cite{b6}, que
avaliam duas técnicas clássicas sob protocolo idêntico e obtêm sinais opostos:
enquanto o \textit{Non-Local Means} eleva o erro médio, a equalização adaptativa 
de histograma (CLAHE) o reduz. A distinção que os
autores propõem é qualitativa: o CLAHE redistribui intensidades,
ao passo que o \textit{denoising} descarta detalhe fino. O critério que separa
melhoria de degradação, portanto, não é alterar a imagem, mas remover
informação.

As demais condições de melhoria confirmam esse critério. O \textit{denoising} beneficia 
a classificação quando a entrada está genuinamente degradada e a dose é 
controlada \cite{b26}, \cite{b8}. A imagem filtrada agrega desempenho quando usada como 
canal adicional ao lado da original \cite{b9}, \cite{b16}. O borramento
moderado funciona como aumento de dados e chega a melhorar a acurácia sobre
imagens nítidas \cite{b4}. A manipulação espectral usada como augmentação de treino, e 
não como filtro de teste, melhora simultaneamente o desempenho limpo e o degradado 
\cite{b28}, \cite{b29}. A transformação de entrada é tolerável quando descarta
seletivamente \cite{b11}. Há, por fim, a condição identificada por Ortiz-Jiménez 
\textit{et al.} \cite{b23}: quando o conjunto de dados simplesmente não contém 
informação discriminativa na banda removida, a filtragem não custa nada.

O padrão que emerge é coerente. O pré-processamento é benigno
quando redistribui em vez de remover, quando é aplicado simetricamente a treino
e teste, quando preserva a imagem original como canal complementar, ou quando
atua sobre bandas que o conjunto de dados não usa. É danoso quando remove,
assimetricamente, energia de bandas nas quais o classificador aprendeu a
decidir.

\subsection{Respostas às Questões de Pesquisa}
\label{sec:rq}

A Tabela~\ref{tab:mapaqp} reporta, para os 20 estudos com extração formal, a
quantas questões de pesquisa cada um foi declarado relevante. A QP3 é a menos
coberta do conjunto --- 13 dos 20 registros, contra 18 para cada uma das
demais ---, o que é consistente com a escassez de evidência quantitativa
discutida adiante.

\begin{table}[htbp]
\caption{Relevância declarada na extração de dados, por questão de pesquisa ($n=20$)}
\begin{center}
\footnotesize
\begin{tabular}{|c|c|c|}
\hline
\textbf{QP} & \textbf{$n$} & \textbf{\% dos 20} \\
\hline
QP1 --- degradação existe & 18 & 90\% \\
\hline
QP2 --- mecanismo explicativo & 18 & 90\% \\
\hline
QP4 --- condições de melhoria & 18 & 90\% \\
\hline
QP3 --- relação quantitativa energia--acurácia & 13 & 65\% \\
\hline
\end{tabular}
\label{tab:mapaqp}
\end{center}
\end{table}

\subsubsection{QP1 --- As técnicas degradam a acurácia?}

\textbf{Sim, com magnitude altamente variável e condicionada ao protocolo.} A
evidência é consistente entre os estudos com \textit{baseline} explícito
reunidos na Tabela~\ref{tab:magnitudes}. A degradação é severa e monotônica
para suavização (90,88\%~$\rightarrow$~60,97\% em \textit{top}-5 sob desfoque
de raio~8 \cite{b4}; queda replicada em seis arquiteturas \cite{b5} e, com
\textit{baseline} limpo, de 84,72\% para 67,84\% no CIFAR-10 \cite{b28}),
substancial para \textit{denoising} aplicado a imagens já limpas
(76,7\%~$\rightarrow$~58,3\% \cite{b7}; mCE de 76,7\%~$\rightarrow$~82,1\%
\cite{b6}), e condicionada a limiar para compressão, que é praticamente inócua
acima de um fator de qualidade moderado \cite{b1}, \cite{b10}.

\subsubsection{QP2 --- Quais mecanismos explicam a degradação?}

\textbf{Quatro mecanismos são propostos na literatura, complementares e não
excludentes.}

\emph{(i) Remoção de sinal discriminativo de alta frequência.} É o mecanismo
mais diretamente apoiado \cite{b14}, \cite{b15}, \cite{b17}.

\emph{(ii) Perda de textura.} Estabelecido por intervenção causal \cite{b16}, \cite{b1}, \cite{b2}. No nível do filtro, afeta preferencialmente os seletivos em frequência e orientação \cite{b25}.

\emph{(iii) Descasamento de domínio.} A
assimetria de \cite{b23} mostra que a maior parte
da degradação observada nos protocolos usuais não decorre de perda
irrecuperável de informação, mas de a entrada passar a ocupar regiões do espaço
onde o modelo tem fronteiras de decisão (corroborado por \cite{b4}).

\emph{(iv) Amplificação de erro na hierarquia.} O resíduo deixado pelo
pré-processamento, pequeno na entrada, cresce progressivamente camada a camada
até dominar as características de topo \cite{b7}. Há ainda uma
causa estrutural subjacente nas funções de base de Fourier \cite{b19}.

\subsubsection{QP3 --- Existe relação quantificável entre energia espectral
removida e queda de acurácia?}

\textbf{Parcialmente, e a relação não é universal.} Existe uma única evidência
de relação quantitativa explícita: Yin \textit{et al.} \cite{b18}. Relações
dose--resposta adicionais existem em forma ordinal \cite{b20} ou restritas a um domínio específico \cite{b21}.

A conclusão operacional é que a energia espectral removida é um preditor
parcial e mediado: mediado pela distribuição espectral do conjunto de treino
\cite{b16}, \cite{b22} e pela simetria do protocolo entre treino e teste
\cite{b23}. Não existe, na literatura revisada, uma curva dose--resposta definitiva construída sobre
técnicas de aprimoramento perceptual e classificadores em modo de inferência.

\subsubsection{QP4 --- Em quais condições o pré-processamento melhora?}

\textbf{Em seis condições bem delimitadas}, sintetizadas na
Tabela~\ref{tab:condicoes}. O critério unificador é que o pré-processamento é benigno quando \emph{redistribui} informação e danoso quando a \emph{descarta}; e que, quando
o descarte é inevitável, seu custo é proporcional à relevância discriminativa
da banda afetada e à assimetria entre treino e teste.

\begin{table}[htbp]
\caption{Condições sob as quais o pré-processamento não degrada ou melhora}
\begin{center}
\footnotesize
\begin{tabular}{|p{3.0cm}|p{1.1cm}|p{2.5cm}|}
\hline
\textbf{Condição} & \textbf{Fonte} & \textbf{Evidência} \\
\hline
A técnica redistribui intensidades em vez de descartar detalhe & \cite{b6} &
CLAHE: erro limpo 23,87\%~$\rightarrow$~23,55\%; mCE
76,7\%~$\rightarrow$~74,5\% \\
\hline
A entrada está genuinamente degradada, o método é casado ao tipo de ruído e a
dose é controlada & \cite{b26}, \cite{b8} & VGG16:
94,6\%~$\rightarrow$~97,0\%; InceptionV3: 84,7\%~$\rightarrow$~90,0\% \\
\hline
A imagem filtrada é usada como canal adicional, não como substituta &
\cite{b9}, \cite{b16} & Canal duplo melhora em todas as degradações;
SIN$+$IN supera a linha de base \\
\hline
A manipulação espectral é usada como augmentação de treino, não como filtro de
teste & \cite{b4}, \cite{b28}, \cite{b29} & Blur moderado:
90,88\%~$\rightarrow$~91,36\% em imagens nítidas; erro de corrupção
28,35\%~$\rightarrow$~26,26\% \\
\hline
O descarte é seletivo e adaptado à imagem & \cite{b11} & Limiar
\textit{wavelet} adaptativo: custo $<$~2 p.p. em imagens limpas \\
\hline
O conjunto de dados não contém sinal discriminativo na banda removida &
\cite{b23} & MNIST sob passa-baixa $14\times14$: perda nula \\
\hline
\end{tabular}
\label{tab:condicoes}
\end{center}
\end{table}

\subsubsection{Lacunas identificadas e posicionamento deste trabalho}

A síntese permite identificar quatro lacunas que motivam a presente investigação: lacuna de \emph{escopo de técnica} \cite{b6}, \cite{b7}, \cite{b8}, \cite{b9}, \cite{b26}, \cite{b3}; lacuna de \emph{variável independente} \cite{b1}, \cite{b4}, \cite{b5}, \cite{b28}, \cite{b13}, \cite{b12}; lacuna de \emph{confundimento entre perda de informação e descasamento de domínio} \cite{b23}; e lacuna de \emph{fragilidade da própria hipótese espectral}. É nessa interseção que o presente trabalho se situa.

%=======================================================================
\begin{thebibliography}{00}

\bibitem{b1} S. Dodge and L. Karam, ``Understanding how image quality affects
deep neural networks,'' in Proc. 8th Int. Conf. Quality of Multimedia
Experience (QoMEX), 2016, pp. 1--6.

\bibitem{b2} S. Karahan, M. K. Yildirim, K. Kirtac, F. S. Rende, G. Butun, and
H. K. Ekenel, ``How image degradations affect deep CNN-based face
recognition?,'' in Proc. Int. Conf. Biometrics Special Interest Group
(BIOSIG), 2016, pp. 1--5.

\bibitem{b3} Y. Blau and T. Michaeli, ``The perception-distortion tradeoff,''
in Proc. IEEE Conf. Computer Vision and Pattern Recognition (CVPR), 2018, pp.
6228--6237.

\bibitem{b4} I. Vasiljevic, A. Chakrabarti, and G. Shakhnarovich, ``Examining
the impact of blur on recognition by convolutional networks,'' arXiv preprint
arXiv:1611.05760, 2017.

\bibitem{b5} P. Roy, S. Ghosh, S. Bhattacharya, and U. Pal, ``Effects of
degradations on deep neural network architectures,'' arXiv preprint
arXiv:1807.10108, 2018.

\bibitem{b6} D. Hendrycks and T. Dietterich, ``Benchmarking neural network
robustness to common corruptions and perturbations,'' in Proc. Int. Conf.
Learning Representations (ICLR), 2019.

\bibitem{b7} F. Liao, M. Liang, Y. Dong, T. Pang, X. Hu, and J. Zhu, ``Defense
against adversarial attacks using high-level representation guided denoiser,''
in Proc. IEEE/CVF Conf. Computer Vision and Pattern Recognition (CVPR), 2018,
pp. 1778--1787.

\bibitem{b8} G. B. Paranhos da Costa, W. A. Contato, T. S. Nazare, J. E. S.
Batista Neto, and M. Ponti, ``An empirical study on the effects of different
types of noise in image classification tasks,'' arXiv preprint
arXiv:1609.02781, 2016.

\bibitem{b9} J. Yim and K.-A. Sohn, ``Enhancing the performance of
convolutional neural networks on quality degraded datasets,'' arXiv preprint
arXiv:1710.06805, 2017.

\bibitem{b10} G. K. Dziugaite, Z. Ghahramani, and D. M. Roy, ``A study of the
effect of JPG compression on adversarial images,'' arXiv preprint
arXiv:1608.00853, 2016.

\bibitem{b11} A. Prakash, N. Moran, S. Garber, A. DiLillo, and J. Storer,
``Deflecting adversarial attacks with pixel deflection,'' in Proc. IEEE/CVF
Conf. Computer Vision and Pattern Recognition (CVPR), 2018, pp. 8571--8580.

\bibitem{b12} H. Chen and K. Ma, ``Enhancing vision: Harmonizing frequency for
imaging quality and perception accuracy,'' in Proc. IEEE Int. Conf. Acoustics,
Speech and Signal Processing (ICASSP), 2025.

\bibitem{b13} Y. Hao, H. Pei, Y. Lyu, Z. Yuan, J.-R. Rizzo, Y. Wang, and Y.
Fang, ``Understanding the impact of image quality and distance of objects to
object detection performance,'' arXiv preprint arXiv:2209.08237, 2022.

\bibitem{b14} J. Jo and Y. Bengio, ``Measuring the tendency of CNNs to learn
surface statistical regularities,'' arXiv preprint arXiv:1711.11561, 2017.

\bibitem{b15} H. Wang, X. Wu, Z. Huang, and E. P. Xing, ``High-frequency
component helps explain the generalization of convolutional neural networks,''
in Proc. IEEE/CVF Conf. Computer Vision and Pattern Recognition (CVPR), 2020.

\bibitem{b16} R. Geirhos, P. Rubisch, C. Michaelis, M. Bethge, F. A. Wichmann,
and W. Brendel, ``ImageNet-trained CNNs are biased towards texture; increasing
shape bias improves accuracy and robustness,'' in Proc. Int. Conf. Learning
Representations (ICLR), 2019.

\bibitem{b17} A. Ilyas, S. Santurkar, D. Tsipras, L. Engstrom, B. Tran, and A.
Madry, ``Adversarial examples are not bugs, they are features,'' in Advances in
Neural Information Processing Systems (NeurIPS), 2019.

\bibitem{b18} D. Yin, R. Gontijo Lopes, J. Shlens, E. D. Cubuk, and J. Gilmer,
``A Fourier perspective on model robustness in computer vision,'' in Advances
in Neural Information Processing Systems (NeurIPS), 2019.

\bibitem{b19} Y. Tsuzuku and I. Sato, ``On the structural sensitivity of deep
convolutional networks to the directions of Fourier basis functions,'' in Proc.
IEEE/CVF Conf. Computer Vision and Pattern Recognition (CVPR), 2019, pp.
51--60.

\bibitem{b20} T. Saikia, C. Schmid, and T. Brox, ``Improving robustness against
common corruptions with frequency biased models,'' in Proc. IEEE/CVF Int. Conf.
Computer Vision (ICCV), 2021, pp. 10211--10220.

\bibitem{b21} J. Frank, T. Eisenhofer, L. Sch\"onherr, A. Fischer, D. Kolossa,
and T. Holz, ``Leveraging frequency analysis for deep fake image recognition,''
in Proc. Int. Conf. Machine Learning (ICML), 2020.

\bibitem{b22} S. R. Maiya, M. Ehrlich, V. Agarwal, S.-N. Lim, T. Goldstein, and
A. Shrivastava, ``A frequency perspective of adversarial robustness,'' arXiv
preprint arXiv:2111.00861, 2021.

\bibitem{b23} G. Ortiz-Jimenez, A. Modas, S.-M. Moosavi-Dezfooli, and P.
Frossard, ``Hold me tight! Influence of discriminative features on deep network
boundaries,'' in Advances in Neural Information Processing Systems (NeurIPS),
2020.

\bibitem{b24} Z. Shao, T. Ren, C. Wang, L. Isik, and D. M. Beck,
``Dissociating spatial frequency reliance from adversarial robustness
advantages in neurally guided deep convolutional neural networks,'' manuscrito
não publicado, s.d.

\bibitem{b25} T. S. Borkar and L. J. Karam, ``DeepCorrect: Correcting DNN
models against image distortions,'' IEEE Trans. Image Processing, vol. 28, no.
12, pp. 6022--6034, 2019.

\bibitem{b26} M. Hami and M. J. JameBozorg, ``Assessing the impact of CNN auto
encoder-based image denoising on image classification tasks,'' arXiv preprint
arXiv:2404.10664, 2024.

\bibitem{b27} N. Entezari and E. E. Papalexakis, ``TensorShield: Tensor-based
defense against adversarial attacks on images,'' arXiv preprint
arXiv:2002.10252, 2020.

\bibitem{b28} J. Wu, ``FCR: Robust image classification method with frequency
domain consistency regularization,'' in Proc. 6th Int. Symp. Computer
Technology and Information Science (ISCTIS), IEEE, 2026.

\bibitem{b29} H. Hwang, K. Lee, and H.-J. Lee, ``Improving corruption
robustness with random erasing in the frequency domain,'' in Proc. Int. Conf.
Electronics, Information, and Communication (ICEIC), IEEE, 2023.

\end{thebibliography}

\end{document}